\begin{frame}{그림으로: 동심원이 3차원에서 두 고도로 벌어짐}
  \begin{center}
  \includegraphics[height=5.3cm]{ch05_3_circles_3d_lift.png}
  \end{center}
  \vskip0.2em
  \footnotesize
  \begin{itemize}
    \item \textbf{왼쪽}: $\phi(x)=(x_1,x_2,x_1^2+x_2^2)$로 올린 3차원 산점도
      --- 바깥쪽 원(파랑)은 $z\approx1.0$, 안쪽 원(빨강)은 $z\approx0.16$.
    \item \textbf{오른쪽}: 그 사이 $z=0.5$ 근방의 \textbf{평평한 분리 평면}
      하나로 두 클래스가 완전히 갈린다.
    \item 3차원의 평면 $\equiv$ 원래 2차원의 원형 경계.
  \end{itemize}
\end{frame}
